%%% FIELD proof-small-size-frontier
By \(\mathrm{ass:k\mbox{-}domain}\), the admissible budgets at a fixed size \(n\) are the integers \(0\leq k\leq (n-1)(n-2)/2\).  We first record a monotonicity consequence of \(\mathrm{def:graph\mbox{-}class}\).  If \(0\leq j\leq k\), then every triple satisfying \(|D\setminus T|\leq j\) also satisfies \(|D\setminus T|\leq k\), while all other membership conditions are unchanged.  Hence
\[
 \mathcal G_{n,j}\subseteq\mathcal G_{n,k}.
\tag{1}
\]
Because \(D_k(n)\) is the maximum of the finite queried-coordinate degrees over \(\mathcal G_{n,k}\), (1) implies
\[
 D_j(n)\leq D_k(n)
\tag{2}
\]
whenever both budgets are admissible.

For \(n=2\), the budget condition gives \(k=0\), and \(\mathrm{lem:tree\mbox{-}maxima\mbox{-}support}\) gives
\[
 D_0(2)=1.
\tag{3}
\]
For \(n=3\), the admissible budgets are \(k=0,1\).  The two-block conclusion of \(\mathrm{lem:small\mbox{-}block\mbox{-}orientation\mbox{-}maxima}\) gives \(d_e(G)\leq1\) for every finite query in either class, and therefore \(D_k(3)\leq1\).  On the other hand, \(\mathrm{lem:tree\mbox{-}maxima\mbox{-}support}\) gives \(D_0(3)=1\); (2) then gives \(D_k(3)\geq1\).  Thus
\[
 D_k(3)=1\qquad(0\leq k\leq1).
\tag{4}
\]

For \(n=4\), the admissible budgets are \(0\leq k\leq3\).  The three-block conclusion of \(\mathrm{lem:small\mbox{-}block\mbox{-}orientation\mbox{-}maxima}\) gives \(D_k(4)\leq2\).  The tree equality \(D_0(4)=2\) from \(\mathrm{lem:tree\mbox{-}maxima\mbox{-}support}\), followed by (2), gives \(D_k(4)\geq2\).  Consequently
\[
 D_k(4)=2\qquad(0\leq k\leq3).
\tag{5}
\]

For \(n=5\), the admissible budgets are \(0\leq k\leq6\).  The tree equality in \(\mathrm{lem:tree\mbox{-}maxima\mbox{-}support}\) gives
\[
 D_0(5)=2.
\tag{6}
\]
When \(k=1\), the refined five-vertex conclusion of \(\mathrm{lem:small\mbox{-}block\mbox{-}orientation\mbox{-}maxima}\) gives \(D_1(5)\leq3\).  The graph-query furnished by \(\mathrm{lem:five\mbox{-}node\mbox{-}cubic\mbox{-}coordinate\mbox{-}transfer}\) belongs to \(\mathcal G_{5,1}\) and has a queried original coefficient of degree three, so the definition of \(D_1(5)\) gives the reverse inequality.  Therefore
\[
 D_1(5)=3.
\tag{7}
\]
For every admissible \(k\geq2\), the four-block conclusion of \(\mathrm{lem:small\mbox{-}block\mbox{-}orientation\mbox{-}maxima}\) gives \(D_k(5)\leq4\).  By \(\mathrm{lem:five\mbox{-}node\mbox{-}quartic\mbox{-}witness}\), there is a graph-query in \(\mathcal G_{5,2}\) whose queried original coefficient has degree four.  Inclusion (1) with \(j=2\) places the same query in \(\mathcal G_{5,k}\), whence \(D_k(5)\geq4\).  Thus
\[
 D_k(5)=4\qquad(2\leq k\leq6).
\tag{8}
\]
Equations (3)--(8) are precisely the displayed table.  Finally, the degree-two tree construction in \(\mathrm{lem:tree\mbox{-}maxima\mbox{-}support}\), the cubic construction in \(\mathrm{lem:five\mbox{-}node\mbox{-}cubic\mbox{-}coordinate\mbox{-}transfer}\), and the quartic construction in \(\mathrm{lem:five\mbox{-}node\mbox{-}quartic\mbox{-}witness}\) each use an original directed coefficient and assert that its distinct generic conjugate values separate the corresponding coordinate branches.  These are exactly the nontrivial lower-bound constructions used above.
%%% FIELD proof-frontier-lower-envelope
Fix an admissible pair \((n,k)\) with \(n\geq4\), and define the auxiliary integer
\[
 j:=\min\!\left\{k,\left\lfloor\frac{n-4}{2}\right\rfloor\right\}.
\tag{9}
\]
Since \(k\geq0\) by \(\mathrm{ass:k\mbox{-}domain}\) and \(n\geq4\), we have \(j\geq0\).  Moreover, \(j\leq\lfloor(n-4)/2\rfloor\), so
\[
 n\geq2j+4.
\tag{10}
\]
Thus \(\mathrm{thm:large\mbox{-}size\mbox{-}frontier}\), applied with excess budget \(j\), gives the one-sink graph-query \((G^{\star}_{j,n},T,e^{\star})\) and the equality
\[
 d_{e^{\star}}(G^{\star}_{j,n})=2^{j+1}.
\tag{11}
\]
Its membership in \(\mathcal G_{n,j}\) is also the explicit conclusion of \(\mathrm{def:attaining\mbox{-}membership\mbox{-}check}\).  Since \(j\leq k\), \(\mathrm{def:graph\mbox{-}class}\) gives
\[
 \mathcal G_{n,j}\subseteq\mathcal G_{n,k}:
\tag{12}
\]
indeed, the sole budget inequality strengthens from \(|D\setminus T|\leq j\) to \(|D\setminus T|\leq k\), and the remaining membership conditions do not change.  The finite degree in (11) is therefore among the values maximized in the definition of \(D_k(n)\).  Equations (11)--(12) imply
\[
 D_k(n)\geq2^{j+1}.
\tag{13}
\]

It remains only to rewrite the exponent.  For every integer \(n\geq4\),
\[
 \left\lfloor\frac{n-4}{2}\right\rfloor+1
 =\left\lfloor\frac{n-2}{2}\right\rfloor.
\tag{14}
\]
For nonnegative integers \(u,v\), one also has \(\min\{u,v\}+1=\min\{u+1,v+1\}\).  Applying this identity to (9) and then using (14) yields
\[
 j+1
 =\min\!\left\{k+1,\left\lfloor\frac{n-2}{2}\right\rfloor\right\}.
\tag{15}
\]
Substitution of (15) into (13) proves
\[
 D_k(n)\geq
 2^{\min\{k+1,\lfloor(n-2)/2\rfloor\}}.
\]
The witness is the same original directed coefficient \(e^{\star}\) specified by \(\mathrm{def:attaining\mbox{-}family}\) and certified by \(\mathrm{thm:large\mbox{-}size\mbox{-}frontier}\).  Finally, \(\mathrm{thm:small\mbox{-}size\mbox{-}frontier}\) gives the separate exact five-vertex conclusions \(D_1(5)=3\) and \(D_k(5)=4\) for every admissible \(k\geq2\), as claimed.
